\section{Proofs and Constants for Task-Relevant Motion}
\label{app:task-geometry}
\subsection{Regression, mobility, and near-conservation}
Centering absorbs the affine intercept. Since $\E uu^\top=I$, minimizing over a decoder $D_Y$ gives
\begin{align}
 \E\norm{Y-D_YAu}^2={}&\tr\operatorname{Cov}(Y)-\norm{C_Y}_F^2\\
 &+\norm{D_YA-C_Y^\top}_F^2.\nonumber
\end{align}
Orthogonal projection onto $\row A$ proves Equation~\eqref{eq:general-task-risk}, including rank-deficient $A$. The differential formulas additionally require constant rank; the explicit bounds below use full row rank and a positive singular-value lower bound. Whitening the covariance support is exact, but estimating it from finite images is a separate statistical problem.

For a full-row-rank $A$, differentiating $P_A=A^\dagger A$ gives Equation~\eqref{eq:projector-tangent}. The two summands are adjoint off-diagonal blocks, whence
\begin{equation}
 \norm{\dot P_A}_F=\sqrt2\norm{A^\dagger\dot A(I-P_A)}_F.
\end{equation}
The head Jacobian is $\delta W\mapsto\delta WB$, with Gram action $G\mapsto GB^\top B$. A channel Jacobian has Gram $T_c$, so $\dot B_{c:}=-\eta_h(W^\top G)_{c:}T_c$. These facts prove Equation~\eqref{eq:mobility-split}. Substituting $G=HWB+E$ and $T_c=\tau_cI+\mathcal E_c$, then multiplying by $Q_B$, eliminates the row-space term $\tau_c(U_0B)_{c:}$ and proves Equation~\eqref{eq:square-leakage}.

For linear $V$, $\dot V=-\eta_VW^\top HWV-\eta_VW^\top E$ proves Equation~\eqref{eq:linear-leakage}. Combining its exterior part with the projector identity yields the stated integral near-conservation bound. When $E=0$, the continuous flow is invertible left multiplication on every finite bounded interval. In ordinary discrete gradient descent each step also left multiplies $V$; retaining its rank suffices for the same row space. Finite step size alone cannot produce drift under the exact isotropic gradient and rank preservation.

Differentiating GUPA with respect to $ASA^\top$ gives $G=AD-KAS$. Since $ASA^\top\succeq0$, $\norm K_2\le2$. Splitting $S,D$ into scalar and residual terms proves Equation~\eqref{eq:gupa-residual}. This is conditional on a common clean/augmented moment support, as specified in the main text; it does not discard view-only coordinates.

\subsection{The finite displacement bound}
Put $P=P_0$, $Q=I-P$, and define
\begin{equation}
 F=P(I-P_1)P,\quad E=QP_1Q,\quad K_1=QP_1P.
\end{equation}
The first two matrices are PSD. Equal ranks imply
\begin{equation}
 \tr F=\tr E=d_{01}^2,\qquad \norm{K_1}_F\le d_{01}.
\end{equation}
For example, principal angles give $d_{01}^2=\sum_i\sin^2\theta_i$ and $\norm{K_1}_F^2=\sum_i\sin^2\theta_i\cos^2\theta_i$. Exact block expansion gives
\begin{align}
 R(P_1)-R(P)={}&\tr(FPTP)-\tr(EQTQ)\\
 &-2\tr(K_1PTQ).\nonumber
 \label{eq:exact-leakage}
\end{align}
Subtracting $\lambda_{\min}(T)I$ leaves the risk difference and task cross block unchanged. Each diagonal trace is then in $[0,\operatorname{osc}(T)d_{01}^2]$, so their difference has absolute value at most the upper endpoint. Cauchy--Schwarz bounds the cross term by $2a_0d_{01}$, proving Equation~\eqref{eq:finite-motion-bound}. If $QTP=0$ and the restricted row/normal spectra have separation $\delta$, the same equality gives the one-sided bound $\delta d_{01}^2$.

Along a full-rank path, integrate the projector derivative to obtain
\begin{equation}
 d_{01}\le\int\norm{A^\dagger\dot A(I-P_A)}_Fdt
 \le\frac J\sigma\int(\norm{\dot h}_F+\norm{\dot W}_F)dt.
\end{equation}
This proves the finite update-budget bound. An invertible rescaling of representation coordinates leaves the risk fixed; a small singular value nevertheless makes a parameter-based sufficient bound more conservative. It does not itself prove information loss.

\subsection{The complete two-block competition criterion}
Let $U,N$ be orthonormal bases of the current row and normal spaces. Every projector tangent has the form
\begin{equation}
 D_g=UK_g^\top N^\top+NK_gU^\top.
\end{equation}
For $\beta=Ux+Ny$, $\beta^\top D_g\beta=2y^\top K_gx$. Thus absence of competition for every rank-one task is equivalent to
\begin{equation}
 (y^\top K_hx)(y^\top K_Wx)\ge0\quad\text{for all }x,y.
\end{equation}
For fixed $x$ with $K_hx\ne0$, this requires $K_Wx=\lambda(x)K_hx$ with $\lambda(x)\ge0$: two noncollinear linear functionals admit opposite signs. If $z\in\ker K_h$, apply the same argument to $x+tz$ for every real $t$. A whole line cannot stay on one nonnegative ray unless $K_Wz=0$. Hence $K_W=F_0K_h$ for a linear map $F_0$ on $\operatorname{im}K_h$. Every vector of that image is its eigenvector with a nonnegative eigenvalue. Linearity applied to two independent vectors and their sum makes that eigenvalue constant. Thus $K_W=\lambda K_h$, with $\lambda>0$ if both tangents are nonzero. The zero-block cases and the converse are immediate. Positive proportionality also excludes competition for every PSD multi-output $T$.

Strict competition holds on an explicit task neighborhood. If a unit task $\beta_0$ has opposing contributions with minimum magnitude $m>0$ and $d=\max_g\norm{D_g}_2$, every unit task with $\norm{\beta-\beta_0}\le m/(4d)$ retains both signs with margin at least $m/2$, since
\begin{equation}
 |\beta^\top D_g\beta-\beta_0^\top D_g\beta_0|
 \le2d\norm{\beta-\beta_0}.
\end{equation}
This states which task directions can compete at a fixed state, not the probability that a prescribed label task competes along random training.

\subsection{Budget windows, finite steps, and explicit constants}
For unit negative block directions, $\norm{d\theta/db}\le1$. On the ball in the main text, each $\langle\nabla R,u_g\rangle$ is $L_q$-Lipschitz. Consequently its deviation from its source value at budget $s$ is at most $L_qs$. Integrating a convex combination proves Equation~\eqref{eq:budget-envelope}. For a constant allocation $\alpha$, the sharper form is
\begin{equation}
 |R_\alpha(b)-R_0-b[\alpha q_h+(1-\alpha)q_W]|
 \le\tfrac12L_qb^2.
\end{equation}
Subtracting two such bounds proves the comparison bound. Source slopes separated from zero by $m$ retain their signs with margin $m/2$ over the displayed half-window. For $q_h=-a<0<q_W=c$, both blocks can be active: allocations $\alpha_*/2$ and $(1+\alpha_*)/2$, where $\alpha_*=c/(a+c)$, have opposite slopes with margin at least $\min(a,c)/2$.

For perturbed Euler updates $\theta_{j+1}=\theta_j+\delta_j(F_\alpha(\theta_j)+e_j)$, $\norm{e_j}\le\varepsilon$, and $b=\sum_j\delta_j$, assume $(1+\varepsilon)b\le\rho$. Taylor's theorem gives
\begin{align}
 |R_N-R_0-bq_\alpha|\le{}&\tfrac12L_q(1+\varepsilon)b^2
 +G_R\varepsilon b\\
 &+\tfrac12H_R(1+\varepsilon)^2\sum_j\delta_j^2.\nonumber
\end{align}
The three terms bound state drift, direction error, and step curvature. Using $\sum\delta_j^2\le\delta_{\max}b$ proves Equation~\eqref{eq:margin-budget} and its cumulative window. The perturbed process uses a nominal budget; its actual path length may differ. If raw block-gradient error is $\delta_g$ and its true norm is at least $\mu_g>0$, normalized-direction error is at most $2\delta_g/\mu_g$. Thus a single batch does not automatically establish a useful uniform $\varepsilon$.

Here are fully explicit sufficient constants for the local square map. Let $k_\phi=\norm{U\Lambda^{1/2}}_2$, and on a radius-$\rho$ parameter ball take $H_f=\norm{h_0}_F+\rho$ and $W_f=\norm{W_0}_2+\rho$. Valid bounds are
\begin{align}
 J&=k_\phi\sqrt{H_f^4+4W_f^2H_f^2},\\
 H_A&=2k_\phi(H_f+W_f),\qquad a=k_\phi W_fH_f^2.
\end{align}
They bound $DA$, $D^2A$, and $\norm A_F$, respectively. A sufficient rank bound is $\sigma=\sigma_{\min}(A_0)-J\rho>0$. One may use
\begin{align}
 G_R&=2\norm T_FJ/\sigma,\\
 H_R&=6\norm T_FJ^2/\sigma^2+2\norm T_FH_A/\sigma.
\end{align}
For verification, $\nabla_AR=-2(A^\dagger)^\top T(I-P_A)$,
\begin{equation}
 DA^\dagger[Z]=-A^\dagger ZA^\dagger
 +(I-P_A)Z^\top(AA^\top)^{-1},
\end{equation}
so $\norm{DA^\dagger[Z]}_2\le2\norm Z_F/\sigma^2$ and $\norm{DP_A[Z]}_2\le\norm Z_F/\sigma$. Differentiation and composition give the displayed risk bounds.

For GUPA put $s_f=\norm S_2$, $d_f=\norm D_2$, $v_f=d_f+2s_f$. The inverse resolvent bound yields
\begin{equation}
 \norm{DG[Z]}_F\le(v_f+6s_f^2a^2)\norm Z_F.
\end{equation}
Hence $H_L=(v_f+6s_f^2a^2)J^2+av_fH_A$ bounds the parameter-loss Hessian. If $\mu_g=\norm{g_g(\theta_0)}-H_L\rho>0$ for both blocks, take $K=H_L/\min_g\mu_g$. These constants require source moments and a specified region, rather than future fitting; they can be very conservative. They have not been shown useful for natural images.

For constant physical rates, ordinary gradient flow has budget allocation $\alpha(\theta)=\eta_h\norm{g_h}/(\eta_h\norm{g_h}+\eta_W\norm{g_W})$, which generally varies with state. Alternatively use clock $d\tau_{\rm rate}=(\eta_h+\eta_W)dt$ and raw negative block fields. If their norms are bounded by $M$ and Lipschitz constants by $K_f$, a corresponding Taylor remainder is $\tfrac12(H_RM+G_RK_f)M\tau_{\rm rate}^2$. A sum of rates is not a parameter path-length budget.

\subsection{Gaussian patches and the two-filter motion criterion}
For $x\sim N(0,I_p)$, Gaussian fourth moments make $u=\operatorname{svec}(xx^\top-I)/\sqrt2$ isotropic. Write $H_c=h_ch_c^\top$. Then $b_c=\sqrt2\operatorname{svec}(H_c)$ and
\begin{equation}
 J_cv=\sqrt2\operatorname{svec}(h_cv^\top+vh_c^\top),
 \quad T_c[S]=4(H_cS+SH_c).
\end{equation}
For $G=HA$ with symmetric $H$, let $K_0=W^\top HW$. The exact encoder update is
\begin{equation}
 \dot b_c=-4\sqrt2\eta_h\operatorname{svec}
 \left(\sum_j(K_0)_{cj}(H_cH_j+H_jH_c)\right).
\end{equation}
The Jordan product is $(h_c^\top h_j)(h_ch_j^\top+h_jh_c^\top)$. Orthogonal or collinear groups are a closure case: those products remain in the current span. For two independent filters put $\rho_h=h_1^\top h_2$, $n=Q_B\operatorname{svec}(h_1h_2^\top+h_2h_1^\top)$, and $p_B=B^\dagger(1,1)^\top$. Both $n,p_B$ are nonzero, and
\begin{equation}
 \dot BQ_B=-4\sqrt2\eta_h(K_0)_{12}\rho_h(1,1)^\top n^\top.
\end{equation}
Thus, for $\eta_h>0$, the encoder moves exactly when $(K_0)_{12}\rho_h\ne0$. Its projector tangent is $vn^\top+nv^\top$ with $v=-4\sqrt2\eta_h(K_0)_{12}\rho_hp_B$. The unit tasks $(v/\norm v\pm n/\norm n)/\sqrt2$ have opposing risk derivatives of magnitude $\norm v\norm n$. A task exactly along $n$ has zero first derivative and can only gain quadratically at the source. This establishes task relevance without asserting that all nonlinear encoders drift or that drift always hurts a task.

\subsection{Proof of Theorem~\ref{thm:bottleneck-family}}
In whitened input coordinates $q_c=\Sigma_x^{1/2}h_c$, Euclidean $h$ training induces metric $\Sigma_x$ on $q$. On the symmetric family $q_1+q_2$ and $q_1-q_2$ remain along $e_+,e_-$; both head entries remain $w$. With $\kappa=2\Phi'(V)$, direct differentiation gives
\begin{align}
 \dot\lambda_\pm&=-8\eta_h\kappa w^2\nu_\pm\lambda_\pm^2,\\
 \dot w&=-\eta_W\kappa w(\lambda_+^2+\lambda_-^2).
\end{align}
Changing to $\tau$ gives Equation~\eqref{eq:analytic-family}. Since $2\nu_+>\nu_-$, the ratio $\lambda_-/\lambda_+$ strictly decreases. The unique state with $0<\theta<\theta_0$ occurs at
\begin{equation}
 \tau(\theta)=\frac{1/2-\tan\theta}{\epsilon(\nu_+-\nu_-\tan\theta)}
 <\frac1{2\nu_+\epsilon}.
\end{equation}
The sum of the two encoder rows is the diagonal direction in the main text; their difference is a nonzero multiple of $E_\times$. This proves the ranks and Equation~\eqref{eq:analytic-risks}. At the two fixed comparison angles, the improvement and excess over initialization are respectively
\begin{align}
 &(1-z^2)\sin(\theta_0/4)\sin(7\theta_0/4-2\psi),\\
 &(1-z^2)\sin(5\theta_0/6)\sin(2\psi-7\theta_0/6).
\end{align}
Their minima on the task cone are the stated $\Delta_1,\Delta_2$. Independent additive label noise simply adds a common constant floor.

It remains to prove the descending regime and finite physical reachability rather than assume them. Population-centered Gaussian views with covariance-relative noise have $S=sI$, $D=2(s-1)I$, $s=(1+\eta)^2\le36/25$. Thus $\kappa=-2\Gamma_s(sV)$ as in the original construction. Since $\theta_0>11/24$, $\tan\theta_2>3/40$. At $\theta_2$ the closed formulas imply
\begin{align}
 \lambda_+&=\frac{2\epsilon(\nu_+/\tan\theta_2-\nu_-)}{2\nu_+-\nu_-}<33\epsilon,\\
 \lambda_-&<\tfrac52\epsilon,\qquad
 w^2<\tfrac1{100}+\tfrac{5}{2}\frac{(65/2)\epsilon}{19/20}<\tfrac{87}{100}.
\end{align}
The modal values and $w^2$ increase up to this state, so throughout the interval
\begin{equation}
 V=2w^2(\lambda_+^2+\lambda_-^2)
 <2\tfrac{87}{100}\big[(33/100)^2+(1/40)^2\big]<\tfrac15.
\end{equation}
Consequently $sV<1/3$, and monotonicity of $\gamma$ gives
\begin{equation}
 \Gamma_s(sV)\ge1-\tfrac{36}{25}(1-\gamma(1/3))
 =\tfrac{101}{200}>\tfrac12.
\end{equation}
This strict barrier closes a first-exit argument for the descending regime. In particular $\dot\tau>8\eta_hw_0^2$, while $\tau_2<1/(2\nu_+\epsilon)\le10/(19\epsilon)$. Therefore $t_2<5/(76\eta_hw_0^2\epsilon)$. The exact formulas bound all parameters before this state, and nonzero front gradients make the loss strictly decrease, including when the head rate is zero.

The covariance interval is an explicit structured anisotropic family. It preserves the alignment and initialization conditions of the theorem. It does not supply a bound for arbitrary perturbations away from this invariant manifold. Similarly, the rate independence of the ordered subspace path is an exclusion result within this family, not a proof that rate bifurcation persists under symmetry breaking.

As an implementation check, \path{experiments/code/theory/analytic_bottleneck_check.py} evaluates 216 closed-form states, compares SVD risks on 25 tasks per state, and differentiates the Gaussian physical quadratic-moment objective in $h,W$ for 648 state/noise combinations. The maximum risk and parameter-velocity discrepancies are about $1.25\times10^{-15}$ and $1.44\times10^{-15}$. The checks include rotated covariance axes, preserve ranks two/one, and satisfy the uniform gaps. They do not replace the analytic proof or constitute random-initialization evidence. No trajectory is trained or integrated.

\subsection{Origin-only interval windows on stored trajectories}
\label{app:origin-windows}
The executable \path{experiments/code/theory/task_motion_check.py} uses the existing fixed-point interval and Taylor-jet arithmetic, but integrates no new trajectory. At each rational source state it constructs a positive forward parameter box $D$. Strictly positive vector-field components exclude exit through lower faces. If $M_i$ bounds the upper component speed, $0.9\min_i(D_i^{\rm upper}-x_{0,i}^{\rm upper})/M_i$ is a valid conservative box lifetime. Fixed spectral powers are not dynamic coordinates. All rational denominators are bounded away from zero on the box.

Taylor jets on the whole box enclose the Lie second derivative $D(DR\,F)F$ by $[q_-,q_+]$; this includes vector-field acceleration, not just the risk Hessian. If $[d_-,d_+]$ encloses the source derivative, then
\begin{equation}
 d_-t+\tfrac12q_-t^2\le R(t)-R_0
 \le d_+t+\tfrac12q_+t^2.
\end{equation}
The chosen duration is the smaller of the box lifetime and half the ratio of the initial slope margin to the second-derivative bound. Integer endpoint checks verify the entire box inclusion and the half-margin inequality after rounding the candidate time downward. Four predetermined relative box widths, $1/100,1/20,1/10,1/5$, are compared using source constants only; the saved trajectory file is opened after all predictions are computed.

The three scales $c=1/2,1,3/2$ use the primary certificate center from Theorem~\ref{thm:gap}, transformed by $W\mapsto c^{-2}W$, $B\mapsto c^2B$. They share $A_0=(0.02432,0.02432)$ and initial risk $1/2$. Their certified physical windows exceed $0.06629,0.09112,0.03072$, respectively. Each contains saved future records; all five records satisfy the predicted amplitude envelopes. These are correlated deterministic trajectory records, not independent statistical trials.

On the same physical boxes the code encloses $0<v_-\le\dot\ell\le v_+$. A first hit of $\ell_*=1/50$ is guaranteed before the box expires when $\ell_*/v_-\le\tau$. Its time lies in $[\ell_*/v_+,\ell_*/v_-]$; substituting this interval into the preceding risk enclosure gives the common-energy predictions in Table~\ref{tab:origin-windows}. Direct differentiation in the energy clock is more conservative: its certified windows are only about $0.01289,0.01081,0.02101$, and only the last directly covers one saved energy increment. The physical-box hitting-time argument is what certifies the common section for all three.

Mapping the scaled states back to the common original parameters gives rates $(c^{-2},c^4)$. After normalizing their sum to one, the three rate pairs are $(64/65,1/65)$, $(1/2,1/2)$, and $(64/793,729/793)$. At common elapsed time $1/10$, source-only risk-change enclosures are conservatively
\begin{align}
 c=1/2:&\quad[0.004267,0.004917],\nonumber\\
 c=1:&\quad[-0.002753,-0.002162],\\
 c=3/2:&\quad[-0.009328,-0.007343].\nonumber
\end{align}
This is a finite split with the same initialization and rate sum. It does not equalize actual parameter path length. The common-energy comparison independently excludes a pure speed explanation. All interval constants and prediction rules precede reading the saved future values; the analysis itself is retrospective, not a new preregistered experiment. The CSV curves are floating-point numerical evidence, not newly certified integrations. The source-only certificates, rather than those curves, establish the displayed forward intervals. No natural-image or stochastic-gradient window is certified here.
